

\begin{landscape}
\begingroup
\setlength{\tabcolsep}{.1pt} % Espacio entre columnas reducido
\renewcommand{\arraystretch}{1.5} % Altura de las filas ligeramente reducida

% 1. El entorno longtable se reemplaza por tabular.
% 2. La tabla completa está dentro de \resizebox para forzar el ajuste al ancho total (landscape).
\begin{table}[h!]
\centering
\caption{Lenguajes de programación: C\#, Python y R}
\label{tab:lenguajesprogramacion}
\footnotesize 
% <-- Ajusta la tabla al ancho de la página horizontal
    % El comando \small reduce el tamaño de la fuente para ayudar al ajuste
    \small 
    \begin{tabular}{|>{\raggedright\arraybackslash}p{2.5cm}|>{\raggedright\arraybackslash}p{6.7cm}|>{\raggedright\arraybackslash}p{7cm}|>{\raggedright\arraybackslash}p{6.7cm}|}
    \hline
    \textbf{Aspecto} & \textbf{C\#}\cite{MarceloVillalobos2008}\cite{M.Deitel2007} & \textbf{Python}\cite{GonzalezDuque}\cite{Downey2002} & \textbf{R}\cite{Paradis2003} \\
    \hline
    \textbf{Caracterí- sticas} & Lenguaje de programación visual orientado a objetos, tipado, estático y controlado por eventos mediante el uso de un IDE. Desarrollado e integrado para el ecosistema .NET  & Lenguaje multiparadigma ya que soporta una programación orientada a objetos, imperativa y funcional con tipado dinámico y con una sintaxis muy fácil de aprender & Lenguaje especializado en análisis estadístico con un enfoque en manipulación de datos y generación de gráficos. \\
    \hline
    \textbf{Rendimi- ento} & Compilado en MSIL y luego optimizado mediante JIT en .NET, de nivel medio. & Interpretado, de alto nivel, lo que implica menor velocidad en tareas intensivas en CPU. Pero permite una gran portabilidad en diferentes entornos de programación. & Interpretado, de alto nivel, optimizado para operaciones estadísticas y matriciales, pero menos eficiente en tareas generales. \\
    \hline
    \textbf{Escalabi- lidad} & Compatible con JSON a través de bibliotecas siendo ampliamente eficiente en sistemas distribuidos y escalables como en ASP.NET. & Soporte nativo para JSON con la biblioteca json y escalable de forma nativa con diferentes \textit{frameworks} para aplicaciones web y con bases de datos modernas. & Compatible con JSON mediante paquetes y escalable en entornos estadísticos de \textit{big data}.\\
    \hline
    \textbf{Ventaja} & 
    Permite una gran integración y soporte multiplataforma robusta con .NET Core al tener un rendimiento optimizado y un tipado estático lo que reduce errores en tiempo de ejecución.
    & 
    Al tener una curva de aprendizaje suave su sintaxis es muy simple y clara, teniendo un ecosistema extenso y una comunidad activa lo que permite una amplia gama de bibliotecas para ciencia de datos y otras herramientas.
    & 
    Permite una visualización muy especializada y un análisis estadístico muy avanzado, siendo que tiene un gran soporte en sus bibliotecas especializadas en el análisis de datos respaldadas por una gran comunidad académica.
    \\
    \hline
    \textbf{Desventaja} & 
    Menor flexibilidad al tener mucha dependencia del ecosistema .NET Core lo que hace gran dependencia a este ecosistema. Esto también hace que algunas herramientas o bibliotecas puedan tener algún costo, aparte el ecosistema .NET puede ser algo robusto en cuestión de memoria.
    & 
    Tiende a tener un rendimiento limitado debido a una concurrencia y un tipado dinámico al hacer la interpretación de línea por línea, lo que lo hace más lento en aplicaciones intensivas. Haciendo que tenga dependencia de bibliotecas externas para algunas aplicaciones.
    & 
    Tiene una flexibilidad muy baja con otros entornos enfocados al desarrollo web, esto debido a su sintaxis única lo que hace que tenga una curva de aprendizaje más pronunciada.
    \\
    \hline
    \end{tabular}

\end{table}

\endgroup
\end{landscape}
